%% Abstract page
%% =============
%%
%% Content of abstract pages has been put into seperate pages to simplify
%% word counting. Use e.g. the unix command
%%   wc abstract-ger.tex
%% or
%%   wc abstract-eng.tex
%% to get the number of words contained in these files.
\thispagestyle{empty}
%\begin{center}
  \noindent\sffamily\textbf{Abstract}
	\\[1em]\normalfont\normalsize

With the increasing number of nuclear-resonant photons per pulse at x-ray free electron lasers (XFELs) \cite{scientific_opp_XFEL}, achieving higher excitation levels experimentally becomes feasible \cite{Shvydko2023}. This development makes it particularly compelling to investigate light propagation and dynamics beyond the low-excitation regime (LER) and around population inversion. Thus, in this thesis, we examine the time response of resonant nuclear scattering of x-rays beyond the LER. In the LER, it is possible to obtain an analytical solution for light propagation, which becomes not possible beyond this regime. To study dynamics beyond the LER, we employ numerical methods to analyze propagation effects and nuclear dynamics. We implement the method of lines (MOL) to solve the Maxwell-Bloch equations.\\

First, we compare our results with those of Wolff \cite{wolff2023characterization} in the LER using thin targets, where the effects arising from the propagation of light through the target can be neglected. We generalize their results by analyzing the dynamics in thick targets, considering the propagation effects. Subsequently, we extend our investigation beyond the LER to the regime of population inversion and beyond. Here, we discover interesting phenomena, such as time shifts in the minima of the observed coherently scattered light intensity and transitions around population inversions. We propose an experimental signature to detect nonlinear excitations by calculating the relative time shifts as a function of target thickness. Finally, we analyze our model in the non-decaying limit ($\gamma = 0$) and find good agreement with the results from Burnham and Chiao \cite{BURNHAM1969}, concluding that spontaneous decay plays a major role in the creation of transitions around population inversions.


  \vfill
  \vspace{1cm}
  %\begin{minipage}[c][0.48\textheight][b]{0.9\textwidth}
    %\small
    %\textbf{
    %  
    %}\par
    %\vspace{\baselineskip}
    %\normalfont\normalsize
    \noindent\sffamily\textbf{Abstrakt}
	\\[1em]\normalfont\normalsize
Mit der zunehmenden Anzahl von nuklear-resonanten Photonen pro Puls beim XFELs \cite{scientific_opp_XFEL} wird es experimentell möglich, höhere Anregungsniveaus zu erreichen \cite{Shvydko2023}. Diese Entwicklung macht es besonders interessant, die Lichtausbreitung und Dynamik über das Niedriganregungsregime (LER) hinaus sowie um die Populationsinversion zu untersuchen. In dieser Arbeit analysieren wir daher das Zeitverhalten der resonanten nuklearen Streuung von Röntgenstrahlen außerhalb des LER. Im LER ist es möglich, eine analytische Lösung für die Lichtausbreitung zu erhalten, was in höheren Anregungsbereichen nicht mehr möglich ist. Um die Dynamik in diesen höheren Bereichen zu untersuchen, verwenden wir numerische Methoden, um die Ausbreitungseffekte und die Kerndynamik zu analysieren. Wir implementieren die Methode der Linien (MOL), um die Maxwell-Bloch-Gleichungen zu lösen.\\

Zunächst vergleichen wir unsere Ergebnisse mit denen von Wolff \cite{wolff2023characterization} im LER unter Verwendung dünner Targets, bei denen die durch die Lichtausbreitung im Target entstehenden Effekte vernachlässigt werden können. Wir erweitern deren Ergebnisse, indem wir die Dynamik in dicken Targets unter Berücksichtigung der Ausbreitungseffekte analysieren. Anschließend untersuchen wir den Bereich der Populationsinversion und darüber hinaus. Hier entdecken wir interessante Phänomene wie Zeitverschiebungen in den Minima der beobachteten kohärent gestreuten Lichtintensität und Übergänge um die Populationsinversionen. Wir schlagen eine experimentelle Signatur vor, um nichtlineare Anregungen durch Berechnung der relativen Zeitverschiebungen in Abhängigkeit von der Targetdicke zu detektieren. Schließlich analysieren wir unser Modell im nicht zerfallenden Grenzfall ($\gamma = 0$) und finden eine gute Übereinstimmung mit den Ergebnissen von Burnham und Chiao \cite{BURNHAM1969}. Wir kommen zu dem Schluss, dass der spontane Zerfall eine wesentliche Rolle bei der Entstehung der Übergänge um die Populationsinversionen spielt.
    %\input{abstract-eng}
  %\end{minipage}
%\end{center}
